% Options for packages loaded elsewhere
\PassOptionsToPackage{unicode}{hyperref}
\PassOptionsToPackage{hyphens}{url}
%
\documentclass[
]{article}
\title{{\'E}valuation\\
ING1 - MSSH-01 - 2026-2027}
\author{}
\date{\vspace{-2.5em}1h30, sans documents}

\usepackage{amsmath,amssymb}
\usepackage{lmodern}
\usepackage{iftex}
\ifPDFTeX
  \usepackage[T1]{fontenc}
  \usepackage[utf8]{inputenc}
  \usepackage{textcomp} % provide euro and other symbols
\else % if luatex or xetex
  \usepackage{unicode-math}
  \defaultfontfeatures{Scale=MatchLowercase}
  \defaultfontfeatures[\rmfamily]{Ligatures=TeX,Scale=1}
\fi
% Use upquote if available, for straight quotes in verbatim environments
\IfFileExists{upquote.sty}{\usepackage{upquote}}{}
\IfFileExists{microtype.sty}{% use microtype if available
  \usepackage[]{microtype}
  \UseMicrotypeSet[protrusion]{basicmath} % disable protrusion for tt fonts
}{}
\makeatletter
\@ifundefined{KOMAClassName}{% if non-KOMA class
  \IfFileExists{parskip.sty}{%
    \usepackage{parskip}
  }{% else
    \setlength{\parindent}{0pt}
    \setlength{\parskip}{6pt plus 2pt minus 1pt}}
}{% if KOMA class
  \KOMAoptions{parskip=half}}
\makeatother
\usepackage{xcolor}
\IfFileExists{xurl.sty}{\usepackage{xurl}}{} % add URL line breaks if available
\IfFileExists{bookmark.sty}{\usepackage{bookmark}}{\usepackage{hyperref}}
\hypersetup{
  pdftitle={Geo Data Science - UE2 Analyse de Données},
  hidelinks,
  pdfcreator={LaTeX via pandoc}}
\urlstyle{same} % disable monospaced font for URLs
\usepackage[left=3cm,right=3cm,top=1cm,bottom=1.5cm]{geometry}
\usepackage{graphicx}
\makeatletter
\def\maxwidth{\ifdim\Gin@nat@width>\linewidth\linewidth\else\Gin@nat@width\fi}
\def\maxheight{\ifdim\Gin@nat@height>\textheight\textheight\else\Gin@nat@height\fi}
\makeatother
% Scale images if necessary, so that they will not overflow the page
% margins by default, and it is still possible to overwrite the defaults
% using explicit options in \includegraphics[width, height, ...]{}
\setkeys{Gin}{width=\maxwidth,height=\maxheight,keepaspectratio}
% Set default figure placement to htbp
\makeatletter
\def\fps@figure{htbp}
\makeatother
\setlength{\emergencystretch}{3em} % prevent overfull lines
\providecommand{\tightlist}{%
  \setlength{\itemsep}{0pt}\setlength{\parskip}{0pt}}
\setcounter{secnumdepth}{-\maxdimen} % remove section numbering
\ifLuaTeX
  \usepackage{selnolig}  % disable illegal ligatures
\fi

\begin{document}
\maketitle

\textit{Répondez sur la copie dans la limite de l'espace disponible (cadre à ne pas dépasser pour chaque question).}

\medskip

\textbf{Nom, Prénom:}

\vspace{3cm}

\subsection{Question 1 : Histoire de la géographie (4 points)} % Géo - Louisette ; répartition des points à rediscuter

En quoi les approches de Paul Vidal de la Blache, figure de la géographie régionale française, et de Waldo Tobler, figure de la géographie quantitative, illustrent-elles deux manières différentes d'étudier l'espace géographique ?

\framebox[\textwidth]{\rule{0pt}{10cm}}

\vspace{1cm}


\newpage


\subsection{Question 2 : Géographie théorique et quantitative (4 points)} % Géo - Louisette

Pourquoi les géographes cherchent-ils à construire des modèles de l'espace géographique ? À partir d'un exemple abordé en cours, expliquez ce qu'un modèle permet de comprendre et présentez une de ses limites.

\framebox[\textwidth]{\rule{0pt}{10cm}}

%\vspace{1cm}
\medskip


\subsection{Question 3 : Sociologie urbaine (4 points)} % Géo - Louisette

Pourquoi la proximité géographique entre les habitants d'une ville ne signifie-t-elle pas nécessairement une proximité sociale ? Expliquez votre réponse en mobilisant au moins deux notions ou exemples abordés en cours.

\framebox[\textwidth]{\rule{0pt}{10cm}}



\newpage


\subsection{Question 4 : Géographie des transports (2 points)} % Géo - Louisette

Quelle différence peut-on faire entre la mobilité et l'accessibilité ? Illustrez cette distinction par un exemple concret.

\framebox[\textwidth]{\rule{0pt}{10cm}}

\medskip


\subsection{Question 5 : Complexité et durabilité (3 points)} % Complexité - Juste

Une planification territoriale nationale est-elle compatible avec une planification urbaine locale, dans le cadre des Objectifs de Développement Durable ?

\framebox[\textwidth]{\rule{0pt}{10cm}}



\newpage


\subsection{Question 6 : Systèmes urbains complexes (3 points)} % Complexité - Juste

A l'aide de vos connaissances et/ou d'exemples vus en cours, illustrez différents aspects de la complexité urbaine.

\framebox[\textwidth]{\rule{0pt}{10cm}}







\end{document}
